% assets/w06-cascade.svg — 바깥 고리는 위치제어, 안쪽 고리는 속도제어
%
%   bash _tools/tikz2svg.sh assets/src/w06-cascade.tex assets/w06-cascade.svg
%
% 이 그림이 말하는 것
%   제어기가 하나가 아니라 **둘이 겹쳐 있다**. 바깥은 "어디로 갈까"를 정하고,
%   안쪽은 "얼마나 세게 밀까"를 정한다. 바깥이 안쪽에게 **명령으로 주는 것은
%   추력이 아니라 속도**다 — 이것이 캐스케이드(다단) 제어의 전부다.
%
%   바깥 고리  위치 오차 [m]   --(P 제어기)-->  속도 지령 [m/s]
%   안쪽 고리  속도 오차 [m/s] --(PI 제어기)--> 추력 지령 [N]
%
% 왜 나누는가 세 가지
%   1) 속도 상한을 **지령 단계에서** 자를 수 있다 — 배가 과속하지 않는다
%   2) 안쪽이 바깥보다 빨라 외란·모델 오차를 먼저 지운다
%   3) 단위가 달라 이득을 따로 맞출 수 있다 (바깥 1/s, 안쪽 N/(m/s))
%
% 좌표 메모 (mm)
%   주 경로 y=0 · 안쪽 되먹임 y=-26 · 바깥 되먹임 y=-44
%   합산점 x = 0 / 96 · 블록 x = 28 / 64 / 126 / 160 / 194
%   블록 폭 18 (반 9) · 합산점 지름 4.6 (반 2.3) — 라벨은 그 사이에 놓는다

\documentclass[border=6pt]{standalone}
% capstone-style.tex — 이 볼트의 TikZ 그림이 공유하는 스타일
%
% 쓰는 법:  % capstone-style.tex — 이 볼트의 TikZ 그림이 공유하는 스타일
%
% 쓰는 법:  % capstone-style.tex — 이 볼트의 TikZ 그림이 공유하는 스타일
%
% 쓰는 법:  \input{capstone-style}   (assets/src/ 안에서)
% 빌드:     bash _tools/tikz2svg.sh assets/src/w02-node-topic.tex assets/w02-node-topic.svg
%
% 왜 TikZ 인가
%   손으로 SVG 좌표를 쓰면 화살촉·선 굵기·글자 크기가 그림마다 미묘하게 달라진다.
%   그것이 "자료가 어설퍼 보인다"의 정체다. 스타일을 한 번만 정하고 상대좌표로
%   배치하면 좌표를 고쳐도 그림이 무너지지 않는다.
%   수식은 본문과 같은 TeX 글꼴로 나온다.
%
% 한글
%   ko.TeX(DVI 경로)를 쓴다. XeLaTeX 은 TikZ 그래픽을 PDF special 로 내보내는데
%   dvisvgm 이 그것을 받으려면 Ghostscript < 10.01 이 필요하다. 이 PC 에는
%   10.07 이 깔려 있어 그 경로가 막힌다. latex -> DVI -> dvisvgm 이 유일하게
%   동작하는 조합이다.

\usepackage[hangul]{kotex}
\usepackage{tikz}
\usepackage{amsmath}
\usepackage{xcolor}
\usetikzlibrary{arrows.meta, positioning, calc, fit, backgrounds, decorations.pathmorphing}

% ---- 색 --------------------------------------------------------------------
% 강조는 하나만. 나머지는 검정. 흑백 인쇄에서도 읽혀야 한다.
\definecolor{ink}{HTML}{1A1A1A}
\definecolor{soft}{HTML}{666666}
\definecolor{accent}{HTML}{7C3AED}
\definecolor{warn}{HTML}{B91C1C}

% 계층 색 — Simulink 모델의 블록 색과 같은 규칙을 쓴다
\definecolor{guid}{HTML}{1D4ED8}    % 유도
\definecolor{ctrl}{HTML}{EA580C}    % 제어
\definecolor{thr}{HTML}{D97706}     % 추진기
\definecolor{plant}{HTML}{16A34A}   % 운동모델
\definecolor{nav}{HTML}{7C3AED}     % 항법 · 신호처리
\definecolor{pathc}{HTML}{0D9488}   % 계획 경로 · 항적
\definecolor{errc}{HTML}{C2410C}    % 오차
\definecolor{flow}{HTML}{0369A1}    % 조류 · 바람

% ---- 선과 화살촉 -----------------------------------------------------------
\tikzset{
  sig/.style   = {ink, line width=0.5pt, -{Stealth[length=4pt, width=3pt]}},
  wire/.style  = {ink, line width=0.5pt},
  fb/.style    = {ink, line width=0.5pt, densely dashed,
                  -{Stealth[length=4pt, width=3pt]}},
  asig/.style  = {accent, line width=0.5pt, densely dashed,
                  -{Stealth[length=4pt, width=3pt]}},
  % 블록 — 흰 바탕, 직각, 검은 테두리. 색은 테두리에만 준다
  blk/.style   = {draw=ink, line width=0.55pt, fill=white,
                  minimum width=18mm, minimum height=8.5mm, inner sep=2pt, font=\small},
  gblk/.style  = {blk, draw=guid},
  cblk/.style  = {blk, draw=ctrl},
  tblk/.style  = {blk, draw=thr},
  pblk/.style  = {blk, draw=plant},
  nblk/.style  = {blk, draw=nav},
  % 합산점 · 분기점
  sum/.style   = {draw=ink, line width=0.55pt, fill=white, circle,
                  inner sep=0pt, minimum size=4.6mm, font=\scriptsize},
  dot/.style   = {ink, fill=ink, circle, inner sep=0pt, minimum size=1.5mm},
  % 글자
  sname/.style = {font=\small, inner sep=1.5pt},
  note/.style  = {font=\scriptsize, soft, inner sep=1.5pt},
  wnote/.style = {font=\scriptsize, warn, inner sep=1.5pt},
  anote/.style = {font=\scriptsize, accent, inner sep=1.5pt},
  axis/.style  = {ink, line width=0.5pt},
  tick/.style  = {font=\scriptsize, soft},
  % 묶음 상자 — 서브시스템 표시
  grp/.style   = {draw=soft, line width=0.4pt, densely dotted, rounded corners=1mm,
                  inner sep=3mm},
}
   (assets/src/ 안에서)
% 빌드:     bash _tools/tikz2svg.sh assets/src/w02-node-topic.tex assets/w02-node-topic.svg
%
% 왜 TikZ 인가
%   손으로 SVG 좌표를 쓰면 화살촉·선 굵기·글자 크기가 그림마다 미묘하게 달라진다.
%   그것이 "자료가 어설퍼 보인다"의 정체다. 스타일을 한 번만 정하고 상대좌표로
%   배치하면 좌표를 고쳐도 그림이 무너지지 않는다.
%   수식은 본문과 같은 TeX 글꼴로 나온다.
%
% 한글
%   ko.TeX(DVI 경로)를 쓴다. XeLaTeX 은 TikZ 그래픽을 PDF special 로 내보내는데
%   dvisvgm 이 그것을 받으려면 Ghostscript < 10.01 이 필요하다. 이 PC 에는
%   10.07 이 깔려 있어 그 경로가 막힌다. latex -> DVI -> dvisvgm 이 유일하게
%   동작하는 조합이다.

\usepackage[hangul]{kotex}
\usepackage{tikz}
\usepackage{amsmath}
\usepackage{xcolor}
\usetikzlibrary{arrows.meta, positioning, calc, fit, backgrounds, decorations.pathmorphing}

% ---- 색 --------------------------------------------------------------------
% 강조는 하나만. 나머지는 검정. 흑백 인쇄에서도 읽혀야 한다.
\definecolor{ink}{HTML}{1A1A1A}
\definecolor{soft}{HTML}{666666}
\definecolor{accent}{HTML}{7C3AED}
\definecolor{warn}{HTML}{B91C1C}

% 계층 색 — Simulink 모델의 블록 색과 같은 규칙을 쓴다
\definecolor{guid}{HTML}{1D4ED8}    % 유도
\definecolor{ctrl}{HTML}{EA580C}    % 제어
\definecolor{thr}{HTML}{D97706}     % 추진기
\definecolor{plant}{HTML}{16A34A}   % 운동모델
\definecolor{nav}{HTML}{7C3AED}     % 항법 · 신호처리
\definecolor{pathc}{HTML}{0D9488}   % 계획 경로 · 항적
\definecolor{errc}{HTML}{C2410C}    % 오차
\definecolor{flow}{HTML}{0369A1}    % 조류 · 바람

% ---- 선과 화살촉 -----------------------------------------------------------
\tikzset{
  sig/.style   = {ink, line width=0.5pt, -{Stealth[length=4pt, width=3pt]}},
  wire/.style  = {ink, line width=0.5pt},
  fb/.style    = {ink, line width=0.5pt, densely dashed,
                  -{Stealth[length=4pt, width=3pt]}},
  asig/.style  = {accent, line width=0.5pt, densely dashed,
                  -{Stealth[length=4pt, width=3pt]}},
  % 블록 — 흰 바탕, 직각, 검은 테두리. 색은 테두리에만 준다
  blk/.style   = {draw=ink, line width=0.55pt, fill=white,
                  minimum width=18mm, minimum height=8.5mm, inner sep=2pt, font=\small},
  gblk/.style  = {blk, draw=guid},
  cblk/.style  = {blk, draw=ctrl},
  tblk/.style  = {blk, draw=thr},
  pblk/.style  = {blk, draw=plant},
  nblk/.style  = {blk, draw=nav},
  % 합산점 · 분기점
  sum/.style   = {draw=ink, line width=0.55pt, fill=white, circle,
                  inner sep=0pt, minimum size=4.6mm, font=\scriptsize},
  dot/.style   = {ink, fill=ink, circle, inner sep=0pt, minimum size=1.5mm},
  % 글자
  sname/.style = {font=\small, inner sep=1.5pt},
  note/.style  = {font=\scriptsize, soft, inner sep=1.5pt},
  wnote/.style = {font=\scriptsize, warn, inner sep=1.5pt},
  anote/.style = {font=\scriptsize, accent, inner sep=1.5pt},
  axis/.style  = {ink, line width=0.5pt},
  tick/.style  = {font=\scriptsize, soft},
  % 묶음 상자 — 서브시스템 표시
  grp/.style   = {draw=soft, line width=0.4pt, densely dotted, rounded corners=1mm,
                  inner sep=3mm},
}
   (assets/src/ 안에서)
% 빌드:     bash _tools/tikz2svg.sh assets/src/w02-node-topic.tex assets/w02-node-topic.svg
%
% 왜 TikZ 인가
%   손으로 SVG 좌표를 쓰면 화살촉·선 굵기·글자 크기가 그림마다 미묘하게 달라진다.
%   그것이 "자료가 어설퍼 보인다"의 정체다. 스타일을 한 번만 정하고 상대좌표로
%   배치하면 좌표를 고쳐도 그림이 무너지지 않는다.
%   수식은 본문과 같은 TeX 글꼴로 나온다.
%
% 한글
%   ko.TeX(DVI 경로)를 쓴다. XeLaTeX 은 TikZ 그래픽을 PDF special 로 내보내는데
%   dvisvgm 이 그것을 받으려면 Ghostscript < 10.01 이 필요하다. 이 PC 에는
%   10.07 이 깔려 있어 그 경로가 막힌다. latex -> DVI -> dvisvgm 이 유일하게
%   동작하는 조합이다.

\usepackage[hangul]{kotex}
\usepackage{tikz}
\usepackage{amsmath}
\usepackage{xcolor}
\usetikzlibrary{arrows.meta, positioning, calc, fit, backgrounds, decorations.pathmorphing}

% ---- 색 --------------------------------------------------------------------
% 강조는 하나만. 나머지는 검정. 흑백 인쇄에서도 읽혀야 한다.
\definecolor{ink}{HTML}{1A1A1A}
\definecolor{soft}{HTML}{666666}
\definecolor{accent}{HTML}{7C3AED}
\definecolor{warn}{HTML}{B91C1C}

% 계층 색 — Simulink 모델의 블록 색과 같은 규칙을 쓴다
\definecolor{guid}{HTML}{1D4ED8}    % 유도
\definecolor{ctrl}{HTML}{EA580C}    % 제어
\definecolor{thr}{HTML}{D97706}     % 추진기
\definecolor{plant}{HTML}{16A34A}   % 운동모델
\definecolor{nav}{HTML}{7C3AED}     % 항법 · 신호처리
\definecolor{pathc}{HTML}{0D9488}   % 계획 경로 · 항적
\definecolor{errc}{HTML}{C2410C}    % 오차
\definecolor{flow}{HTML}{0369A1}    % 조류 · 바람

% ---- 선과 화살촉 -----------------------------------------------------------
\tikzset{
  sig/.style   = {ink, line width=0.5pt, -{Stealth[length=4pt, width=3pt]}},
  wire/.style  = {ink, line width=0.5pt},
  fb/.style    = {ink, line width=0.5pt, densely dashed,
                  -{Stealth[length=4pt, width=3pt]}},
  asig/.style  = {accent, line width=0.5pt, densely dashed,
                  -{Stealth[length=4pt, width=3pt]}},
  % 블록 — 흰 바탕, 직각, 검은 테두리. 색은 테두리에만 준다
  blk/.style   = {draw=ink, line width=0.55pt, fill=white,
                  minimum width=18mm, minimum height=8.5mm, inner sep=2pt, font=\small},
  gblk/.style  = {blk, draw=guid},
  cblk/.style  = {blk, draw=ctrl},
  tblk/.style  = {blk, draw=thr},
  pblk/.style  = {blk, draw=plant},
  nblk/.style  = {blk, draw=nav},
  % 합산점 · 분기점
  sum/.style   = {draw=ink, line width=0.55pt, fill=white, circle,
                  inner sep=0pt, minimum size=4.6mm, font=\scriptsize},
  dot/.style   = {ink, fill=ink, circle, inner sep=0pt, minimum size=1.5mm},
  % 글자
  sname/.style = {font=\small, inner sep=1.5pt},
  note/.style  = {font=\scriptsize, soft, inner sep=1.5pt},
  wnote/.style = {font=\scriptsize, warn, inner sep=1.5pt},
  anote/.style = {font=\scriptsize, accent, inner sep=1.5pt},
  axis/.style  = {ink, line width=0.5pt},
  tick/.style  = {font=\scriptsize, soft},
  % 묶음 상자 — 서브시스템 표시
  grp/.style   = {draw=soft, line width=0.4pt, densely dotted, rounded corners=1mm,
                  inner sep=3mm},
}


\begin{document}
\begin{tikzpicture}[x=1mm, y=1mm]

% ===========================================================================
% 주 경로
% ===========================================================================
\node[sum]  (s1)  at (0,0)   {};
\node[font=\scriptsize] at (0,0) {$+$};
\node[cblk] (kp)  at (28,0)  {위치 P 제어기};
\node[cblk] (sat) at (64,0)  {속도 제한};
\node[sum]  (s2)  at (96,0)  {};
\node[font=\scriptsize] at (96,0) {$+$};
\node[cblk] (pi)  at (126,0) {속도 PI 제어기};
\node[tblk] (al)  at (160,0) {추력 배분};
\node[pblk] (pl)  at (194,0) {배의 운동};

% 합산점 부호 — 지령은 왼쪽(+), 되먹임은 아래(-)
\node[font=\scriptsize, anchor=north east] at (-1.6,-1.8)  {$-$};
\node[font=\scriptsize, anchor=north east] at (94.4,-1.8) {$-$};

% --- 선과 신호 이름 ---------------------------------------------------------
\draw[sig] (-28,0) -- (s1);
\node[note, anchor=south] at (-15,1.4) {$p_d$ \ [m]};
\node[font=\scriptsize, anchor=north] at (-15,-6.2) {목표 위치};

\draw[sig] (s1) -- (kp);
\node[note, anchor=south] at (10.6,1.4) {$e_p$ \ [m]};

\draw[sig] (kp) -- (sat);
\node[anote, anchor=south] at (46,1.4) {$u_d$ \ [m/s]};
\node[font=\scriptsize, accent, anchor=north] at (46,-6.2) {가고 싶은 속도};

\draw[sig] (sat) -- (s2);
\node[anote, anchor=south] at (83.4,1.4) {$u_{d,\,sat}$};

\draw[sig] (s2) -- (pi);
\node[note, anchor=south] at (107.6,1.4) {$e_u$ \ [m/s]};

\draw[sig] (pi) -- (al);
\node[note, anchor=south] at (143,1.4) {$\boldsymbol\tau$};
\node[font=\scriptsize, anchor=north] at (143,-6.2) {[N, N$\cdot$m]};

\draw[sig] (al) -- (pl);
\node[note, anchor=south] at (177,1.4) {$\mathbf{u}$};
\node[font=\scriptsize, anchor=north] at (177,-6.2) {추진기 힘 [N]};

\draw[wire] (pl) -- (232,0);
\node[note, anchor=south] at (218,1.4) {$p,\ u$};

% ===========================================================================
% 되먹임 둘
% ===========================================================================
\node[dot] (d1) at (214,0) {};
\node[dot] (d2) at (226,0) {};

\draw[fb] (d1) |- (96,-26) -- (s2);
\node[anote, anchor=south] at (160,-25) {$u$ \ [m/s] — \textbf{실제 속도}};

\draw[fb] (d2) |- (0,-44) -- (s1);
\node[note, anchor=south] at (160,-43) {$p$ \ [m] — \textbf{실제 위치}};

% ===========================================================================
% 두 고리를 영역으로 묶는다
% ===========================================================================
\draw[grp, draw=ctrl] (86,-32) rectangle (220,12);
\node[font=\scriptsize\bfseries, ctrl, anchor=north west] at (87,-32.8)
     {안쪽 고리 — 속도제어 (빠르다)};

\draw[grp, draw=guid] (-34,-50) rectangle (234,20);
\node[font=\scriptsize\bfseries, guid, anchor=north west] at (-33,-50.8)
     {바깥 고리 — 위치제어 (느리다)};

% ===========================================================================
% 두 고리가 하는 일 — 한 줄씩
% ===========================================================================
\node[font=\scriptsize, guid, anchor=south west, text width=104mm] at (-34,24)
     {\textbf{바깥 고리가 묻는 것} — \emph{"선석까지 12 m 남았다. 그러면 몇 m/s 로 갈까?"} \\
      답은 남은 거리에 비례한다. 멀면 빠르게, 가까우면 느리게.
      그래서 $u_d = K_p\,e_p$ — 나오는 것은 \textbf{속도}이지 힘이 아니다};

\node[font=\scriptsize, ctrl, anchor=south west, text width=116mm] at (118,24)
     {\textbf{안쪽 고리가 묻는 것} — \emph{"1.2 m/s 로 가라는데 지금 0.8 m/s 다. 얼마나 더 밀까?"} \\
      답은 속도 오차에 비례하고, 남은 치우침은 적분이 지운다.
      그래서 $X = K_P\,e_u + K_I\!\int\! e_u\,dt$ — 나오는 것이 \textbf{힘}이다};

% ===========================================================================
% 아래 띠 — 왜 하나로 하지 않고 둘로 나누는가
% ===========================================================================
\draw[wire, soft] (-34,-64) -- (234,-64);

\node[anchor=north west, font=\scriptsize, text width=82mm] at (-34,-69)
     {\textbf{하나로 하면 안 되는가} \\[3pt]
      위치 오차에서 추력을 바로 뽑는 한 고리짜리도 만들 수는 있다. 다만 \\[2pt]
      $\bullet$ 배가 \textbf{얼마나 빨리 가는지를 제어기가 모른다}.
        12 m 남았을 때 추력을 세게 주면 배는 계속 빨라지기만 한다 \\[2pt]
      $\bullet$ 속도를 막으려면 \textbf{추력}을 잘라야 하는데,
        추력을 자르면 멈추는 능력까지 같이 잘린다 \\[2pt]
      $\bullet$ 물 저항과 바람이 바뀌면 같은 추력에 다른 속도가 나온다.
        고칠 자리가 없다};

\node[anchor=north west, font=\scriptsize, text width=82mm] at (58,-69)
     {\textbf{둘로 나누면 생기는 것} \\[3pt]
      $\bullet$ \textcolor{accent}{속도 지령을 자르는 자리}가 생긴다.
        그 값이 곧 \textbf{배가 낼 수 있는 최고 속도}다.
        접안에서는 이것이 안전 장치 노릇을 한다 \\[2pt]
      $\bullet$ 안쪽 고리가 \textbf{바깥보다 빠르게} 돌아, 저항이 바뀌어도
        속도를 지령에 맞춰 놓는다. 바깥 고리는 그 사실을 몰라도 된다 \\[2pt]
      $\bullet$ 이득을 따로 맞춘다. 바깥 $K_p$ 의 단위는 $1/\mathrm{s}$,
        안쪽 $K_P$ 는 $\mathrm{N}/(\mathrm{m/s})$ — 서로 섞이지 않는다};

\node[anchor=north west, font=\scriptsize, text width=82mm] at (150,-69)
     {\textbf{빠르고 느리다는 말의 뜻} \\[3pt]
      안쪽 고리는 \textbf{속도}를 다루고, 속도는 추력에 곧바로 반응한다.
      바깥 고리는 \textbf{위치}를 다루고, 위치는 속도를 한 번 더 쌓아야 바뀐다 \\[2pt]
      그래서 같은 배라도 안쪽이 먼저 자리를 잡고, 바깥은 그 결과를 보고 천천히 움직인다.
      둘의 반응 속도가 비슷해지면 서로 간섭해 흔들린다 \\[2pt]
      $\;\to\;$ 바깥 이득은 안쪽보다 \textbf{작게} 잡는다};

% ===========================================================================
% 6주차 모델에서 이 그림이 어디인가
% ===========================================================================
\node[anchor=north west, font=\scriptsize, text width=266mm] at (-34,-122)
     {\textbf{6주차 모델에서는} — 이 한 쌍이 \textbf{세 벌} 돈다.
      앞뒤 방향($x$)에 한 벌, 좌우 방향($y$)에 한 벌, 선수각($\psi$)에 한 벌이다.
      선수각 고리만 적분 없이 비례-미분으로 둔다 — 각도 오차를 쌓으면 한 바퀴 돌아 버리기 때문이다.
      세 고리가 낸 $X,\ Y,\ N$ 을 \textcolor{thr}{추력 배분}이 받아 추진기 셋의 힘 $F_L,\ F_R,\ F_B$ 로 나눈다};

\end{tikzpicture}
\end{document}
